\documentclass{article}
\usepackage[most]{tcolorbox}
\usepackage{array}
\usepackage[table,dvipsnames]{xcolor}
\usepackage{graphicx, float}
\usepackage[T1]{fontenc}
\usepackage[utf8]{inputenc}
\usepackage[sfdefault]{atkinson}
\graphicspath{{images/}}
\usepackage{amsthm, amssymb, amsmath, bbold}
\usepackage[a4paper]{geometry}
\usepackage[skip=8pt, indent=0pt]{parskip}
\usepackage{pgfplots}
\usepackage[hidelinks]{hyperref}
\usepackage{cancel}
\pgfplotsset{width=8cm,compat=1.9}
\geometry{
a4paper,
total={170mm,257mm},
left=20mm,
top=20mm,
}

\title{Calcolo Numerico}
\author{Marco Pittarello}
\date{}

\newtcbtheorem{theo}{Teorema}{
    colframe=BrickRed!100!white,
    colback=red!5!white
}{th}

\newtcbtheorem{ex}{Esempio}{
    colframe=OliveGreen!100!white,
    colback=OliveGreen!5!white
}{ex}

\newtcbtheorem{defi}{Definizione}{
    colframe=blue!100!white,
    colback=blue!10!white
}{def}

\newtcbtheorem{note}{Nota}{
    colframe=gray!80!black,
    colback=gray!5!white
}{nt}

\renewcommand{\arraystretch}{2}
\setlength{\tabcolsep}{20pt}

\begin{document}
\maketitle
\tableofcontents
\newpage

\section{Sistemi non decimali}

\begin{defi}{Conversione da base $N$ a base $M$}{}
    Dato $x>0$ e nota la rappresentazione $(x)_N$, vogliamo trovare la rappresentazione $(x)_M$.

    Scriviamo:
    \begin{align*}
        x
         & =x_nN^n+x_{n-1}N^{n-1}+\dots+x_0
        +x_{-1}N^{-1}+\dots+x_{-r}N^{-r}    \\
         & =y_mM^m+y_{m-1}M^{m-1}+\dots+y_0
        +y_{-1}M^{-1}+\dots+y_{-q}M^{-q}.
    \end{align*}

    Si procede separatamente per la \underline{parte intera} e per la \underline{parte frazionaria}.
\end{defi}

\subsection{Conversione della parte intera}

Per convertire la parte intera si itera lo schema della divisione intera per $M$:
\begin{align*}
    y_0     & :=x\bmod M
            &
    x       & =q_0M+y_0        \\
    y_1     & :=q_0\bmod M
            &
    q_0     & =q_1M+y_1        \\
            & \vdots           \\
    y_m     & :=q_{m-1}\bmod M
            &
    q_{m-1} & =y_m<M
\end{align*}

\colorbox{yellow}{\textbf{OSSERVAZIONE:}}
$x\bmod M$ indica il \underline{resto della divisione intera} di $x$ per $M$.

I resti ottenuti sono le cifre della rappresentazione in base $M$, lette dal basso verso l'alto. Pertanto:

$$
    (x_{\mathrm{int}})_M
    =
    y_my_{m-1}\dots y_1y_0
$$

\begin{ex}{Conversione da base $2$ a base $10$}{}
    Convertiamo

    $$
        (10011010010)_2
    $$

    in base $10$.

    Si ha:
    \begin{align*}
        (10011010010)_2
         & =
        1\cdot2^{10}
        +0\cdot2^9
        +0\cdot2^8
        +1\cdot2^7
        +1\cdot2^6             \\
         & \qquad
        +0\cdot2^5
        +1\cdot2^4
        +0\cdot2^3
        +0\cdot2^2
        +1\cdot2^1
        +0\cdot2^0             \\
         & =
        2^{10}+2^7+2^6+2^4+2^1 \\
         & =
        1024+128+64+16+2       \\
         & =
        1234
    \end{align*}

    Quindi:

    $$
        \boxed{(10011010010)_2=(1234)_{10}}
    $$

\end{ex}

\begin{ex}{Conversione da base $10$ a base $2$}{}
    Convertiamo

    $$
        (1234)_{10}
    $$

    in base $2$ utilizzando divisioni successive per $2$:
    \begin{align*}
        1234 & =617\cdot2+0 \\
        617  & =308\cdot2+1 \\
        308  & =154\cdot2+0 \\
        154  & =77\cdot2+0  \\
        77   & =38\cdot2+1  \\
        38   & =19\cdot2+0  \\
        19   & =9\cdot2+1   \\
        9    & =4\cdot2+1   \\
        4    & =2\cdot2+0   \\
        2    & =1\cdot2+0   \\
        1    & =0\cdot2+1
    \end{align*}

    Leggendo i resti dal basso verso l'alto otteniamo:

    $$
        10011010010
    $$

    Pertanto:

    $$
        \boxed{(1234)_{10}=(10011010010)_2}
    $$

\end{ex}

\subsection{Conversione della parte frazionaria}

Per convertire la parte frazionaria si procede moltiplicando successivamente per $M$.

Sia

$$
    x_{\mathrm{fraz}}\in[0,1).
$$

Si calcola:

$$
    Mx_{\mathrm{fraz}}=y_{-1}+r_1,
    \qquad
    0\le r_1<1,
$$

dove $y_{-1}$ è la parte intera del risultato.

Si ripete poi il procedimento:

$$
    Mr_1=y_{-2}+r_2,
$$

$$
    Mr_2=y_{-3}+r_3,
$$

e così via.

In generale:

$$
    Mr_k=y_{-(k+1)}+r_{k+1},
    \qquad
    0\le r_{k+1}<1.
$$

Se per qualche $k$ si ottiene

$$
    r_k=0,
$$

il procedimento termina e la rappresentazione frazionaria è finita.

Le cifre ottenute, questa volta, si leggono \underline{dall'alto verso il basso}:

$$
    (x_{\mathrm{fraz}})_M
    =
    (0.y_{-1}y_{-2}y_{-3}\dots)_M.
$$

\begin{ex}{Conversione di una frazione decimale in binario}{}
    Convertiamo

    $$
        (0.0625)_{10}
    $$

    in base $2$.

    Moltiplichiamo successivamente per $2$:
    \begin{align*}
        0.0625\cdot2 & =0.125
                     &        & \longrightarrow 0 \\
        0.1250\cdot2 & =0.250
                     &        & \longrightarrow 0 \\
        0.2500\cdot2 & =0.500
                     &        & \longrightarrow 0 \\
        0.5000\cdot2 & =1.000
                     &        & \longrightarrow 1
    \end{align*}

    La parte frazionaria successiva è $0$, quindi il procedimento termina.

    Pertanto:

    $$
        \boxed{(0.0625)_{10}=(0.0001)_2}
    $$

\end{ex}

\begin{ex}{Un numero frazionario periodico}{}
    Consideriamo il numero

    $$
        (0.1)_{10}.
    $$

    Procedendo come sopra:
    \begin{align*}
        0.1\cdot2 & =0.2 &  & \longrightarrow 0 \\
        0.2\cdot2 & =0.4 &  & \longrightarrow 0 \\
        0.4\cdot2 & =0.8 &  & \longrightarrow 0 \\
        0.8\cdot2 & =1.6 &  & \longrightarrow 1 \\
        0.6\cdot2 & =1.2 &  & \longrightarrow 1 \\
        0.2\cdot2 & =0.4 &  & \longrightarrow 0
    \end{align*}
    e da questo punto la sequenza si ripete.

    Quindi:

    $$
        \boxed{
            (0.1)_{10}
            =
            (0.0001100110011\ldots)_2
        }
    $$

    Il numero ha quindi una rappresentazione binaria infinita e periodica.

    \colorbox{yellow}{\textbf{OSSERVAZIONE:}}
    Un calcolatore binario non può rappresentare esattamente, in generale, un numero come $0.1$ utilizzando un numero finito di bit.
\end{ex}

\section{Numeri macchina}

\subsection{Interi macchina}

I tipi interi macchina più comuni utilizzano $32$ o $64$ bit, ad esempio \texttt{int32} e \texttt{int64}.

\colorbox{yellow}{\textbf{NOTA:}}
Gli interi macchina a $32$ o $64$ bit non fanno parte dello standard IEEE 754, che riguarda invece la rappresentazione dei numeri in virgola mobile.

Se il primo bit è riservato al segno:

$$
    0\longrightarrow +
    \qquad
    1\longrightarrow -
$$

Con $n$ bit e rappresentazione \underline{segno-modulo}, il massimo intero positivo rappresentabile è:

$$
    I_{\max,n}
    =
    (011\dots1)_2
    =
    \sum_{i=0}^{n-2}2^i
    =
    2^{n-1}-1
$$

Per una rappresentazione \underline{senza segno}, invece:

$$
    I_{\max,n}
    =
    (111\dots1)_2
    =
    \sum_{i=0}^{n-1}2^i
    =
    2^n-1
$$

\subsubsection{Overflow}

Si verifica un \underline{overflow} quando il risultato di un'operazione non è rappresentabile con il numero di bit disponibili.

\begin{ex}{Overflow con interi macchina}{}
    Consideriamo interi macchina a $4$ bit in rappresentazione segno-modulo.

    Calcoliamo:

    $$
        7+2
    $$

    Abbiamo:

    $$
        7=(0111)_2
        \qquad
        2=(0010)_2
    $$

    La somma produce:

    $$
        \begin{array}{c}
            0111 \\
            0010 \\
            \hline
            1001
        \end{array}
    $$

    Il risultato ottenuto è:

    $$
        (1001)_2=-1
    $$

    nella rappresentazione segno-modulo.

    Quindi la macchina produce:

    $$
        \boxed{7+2=-1}
    $$

    anziché $9$.

    Questo fenomeno prende il nome di \underline{overflow}.
\end{ex}

\subsection{Floating point e standard IEEE 754}

Usando la notazione binaria, un numero reale $x\neq0$ può essere scritto, in forma normalizzata, come:

$$
    (x)_2
    =
    (-1)^s\cdot 2^{e-b}\cdot(1.f)
$$

dove:

\begin{itemize}
    \item $s$ è il \underline{bit di segno};
    \item $e$ è l'\underline{esponente memorizzato};
    \item $b$ è il \underline{bias}, cioè la deviazione utilizzata per memorizzare anche esponenti negativi;
    \item $e-b$ è l'\underline{esponente effettivo};
    \item $1.f$ è la \underline{significanda} (spesso chiamata informalmente mantissa).
\end{itemize}

Per un numero normalizzato, il bit iniziale $1$ è sempre presente e quindi non è necessario memorizzarlo esplicitamente: si parla di \underline{hidden bit}.

Il valore $0$ richiede invece una convenzione particolare.

\subsection{Tabella riassuntiva IEEE 754}

\begin{center}
    \begin{tabular}{|c|c|c|}
        \hline
                                    & \textbf{32 bit} & \textbf{64 bit} \\
        \hline
        Bit totali                  & 32              & 64              \\
        \hline
        Bit segno                   & 1               & 1               \\
        \hline
        Bit esponente               & 8               & 11              \\
        \hline
        Bit frazione                & 23              & 52              \\
        \hline
        Bias $b$                    & 127             & 1023            \\
        \hline
        Intervallo $e$ normalizzato & $0<e<255$       & $0<e<2047$      \\
        \hline
        $e=0,\ f=0$                 & $\pm0$          & $\pm0$          \\
        \hline
        $e=0,\ f>0$                 & denormalizzato  & denormalizzato  \\
        \hline
        $e=255,\ f=0$               & $\pm\infty$     & $\pm\infty$     \\
        \hline
        $e=255,\ f>0$               & NaN             & NaN             \\
        \hline
    \end{tabular}
\end{center}

\subsection{Arrotondamento e troncamento}

Quando il numero possiede più cifre di quante ne possano essere memorizzate, è necessario decidere cosa fare delle cifre in eccesso.

\begin{defi}{Troncamento}{}
    Con il \underline{troncamento}, le cifre oltre la precisione disponibile vengono semplicemente eliminate.
\end{defi}

\begin{defi}{Arrotondamento}{}
    Con l'\underline{arrotondamento}, si considera la prima cifra che non può essere memorizzata e si modifica l'ultima cifra memorizzabile in funzione del suo valore.

    Nel caso binario:
    \begin{itemize}
        \item se la cifra successiva è $1$, si aumenta di $1$ l'ultima cifra memorizzata;
        \item se la cifra successiva è $0$, l'ultima cifra non viene modificata.
    \end{itemize}
\end{defi}
\newpage
\section{Errore e norme}

\begin{defi}{Approssimazione}{}
    Supponiamo di voler rappresentare o calcolare un valore $x\in\mathbb{R}$ $\left(\text{o }\in\mathbb{R}^m\right)$,
    ma invece calcolo o ottengo un'\underline{approssimazione}: $$\widetilde{x}\approx x$$
\end{defi}

\subsection{Errore assoluto e relativo}

Se $x\in\mathbb{R}$, definiamo l'\underline{errore assoluto} come:

$$
    \boxed{
        |\widetilde{x}-x|
    }
$$

Se, invece, $x\in\mathbb{R}^n$, per rappresentare l'errore assoluto, si utilizza:

$$
    \boxed{
        \|\widetilde{x}-x\|
    }
$$
Dove $||\cdot||$ è una norma su $\mathbb R^m$.

Se $x\in\mathbb{R}\setminus\{0\}$, allora definisco l'\underline{errore relativo} come:$$\boxed{\frac{|\widetilde{x}-x|}{|x|}}$$

\begin{defi}{Norma}{}
    Sia $V$ uno spazio vettoriale e sia:

    $$
        \|\cdot\|:V\longrightarrow\mathbb{R}_{\ge0}
    $$

    La funzione $\|\cdot\|$ è una norma se soddisfa, per ogni $u,v\in V$ e per ogni $\lambda\in\mathbb{R}$:

    \begin{enumerate}
        \item \textbf{omogeneità}: $\|\lambda v\|=|\lambda|\,\|v\|$

        \item \textbf{definitezza}: $\|v\|=0\quad\Longrightarrow\quad v=0_V$

        \item \textbf{disuguaglianza triangolare}: $\|u+v\|\le\|u\|+\|v\|$
    \end{enumerate}
\end{defi}

\subsection{Norma euclidea}

Per

$$
    v=(v_1,\dots,v_n)^T\in\mathbb{R}^n
$$

la norma euclidea è definita da:

$$
    \boxed{
        \|v\|
        =
        \sqrt{v_1^2+v_2^2+\dots+v_n^2}
        =
        \left(
        \sum_{d=1}^{n}v_d^2
        \right)^{1/2}}
$$

\subsection{Norma $p$}

Per $p\ge1$ si definisce la norma $p$ come:

$$
    \boxed{
        \|v\|_p
        =
        \left(
        \sum_{d=1}^{n}|v_d|^p
        \right)^{1/p}
    }
$$
\subsection{Norma infinito}
\begin{defi}{Norma infinito}{}
    La norma infinito è definita come:
    $$
        \|v\|_\infty
        :=
        \max_i |v_i|
    $$
\end{defi}
\subsection{Equivalenza delle norme}

Sia $\{v_k\}$ una successione in $\mathbb{R}^n$ e siano $\|\cdot\|_a,\ \|\cdot\|_b$ due norme.

In $\mathbb{R}^n$ tutte le norme sono equivalenti. In particolare:


\begin{align*}
     & v_k\longrightarrow v
    \quad\text{rispetto a }\|\cdot\|_a\quad\left(\text{ossia } \lim_{k\to+\infty}||v_k-v||_a=0\right) \\
     & \Updownarrow                                                                                   \\
     & v_k\longrightarrow v
    \quad\text{rispetto a }\|\cdot\|_b\quad\left(\text{ossia } \lim_{k\to+\infty}||v_k-v||_b=0\right)
\end{align*}


\subsection{Errore relativo}
Per $v\in\mathbb R^m$ e $\tilde{v}$ appross. di $v$ definiamo l'\underline{errore relativo} rispetto alla norma $\|\cdot\|$ su $\mathbb R^m$ come:

$$
    \boxed{
        \operatorname{Err}_{\text{rel}}(\widetilde{v},\|\cdot\|)
        =
        \frac{\|\widetilde{v}-v\|}{\|v\|}
    }
$$

\subsection{Floating point: successivo e precedente}

Poiché $\mathbb{F}_{32}$ e $\mathbb{F}_{64}$ sono insiemi finiti e ordinati, possiamo definire:

\begin{defi}{Successivo e precedente}{}
    Il \underline{successivo} di $x$:

    $$
        \operatorname{succ}(x)
        =
        \min\{y\in\mathbb{F}:y>x\}\quad \forall x\in\mathbb F \setminus\{\max\mathbb F\}
    $$

    Il \underline{precedente} di $x$:

    $$
        \operatorname{prec}(x)
        =
        \max\{y\in\mathbb{F}:y<x\}\quad \forall x\in\mathbb F \setminus\{\min\mathbb F\}
    $$
\end{defi}

Quanto vale $\operatorname{succ}(x)-x$?

Pongo
\begin{align*}
    x                      & = (-1)^s \cdot 2^{e-b} \cdot 1.f \\
    \operatorname{succ}(x) & =
    \begin{cases}
        (-1)^s \cdot 2^{e-b} \cdot 1.f + (-1)^s \cdot 2^{e-b} \cdot 0.\overbrace{00\dots0}^{m-1}1 & , s=0 \\
        x - (-1)^s \cdot 2^{e-b} \cdot 0.\overbrace{00\dots0}^{m-1}1                              & , s=1
    \end{cases}
\end{align*}

Dato che $f$ contiene $n$ bit, allora la distanza tra due numeri floating point consecutivi, nello stesso intervallo di esponente, è:

$$
    \boxed{
        \operatorname{succ}(x)-x
        =
        2^{e-b-n}
    }
$$

\subsection{Errore di rappresentazione}
$$
    |x-\operatorname{fl}^{T}(x)|
    \le
    \begin{cases}
         & \operatorname{succ}(\widetilde{x})-\widetilde{x}\quad x>0   \\
         & \widetilde{x} - \operatorname{succ}(\widetilde{x})\quad x<0
    \end{cases}
    \ =
    |\operatorname{succ}(\widetilde{x})-\widetilde{x}|
    =
    2^{e-b-n}
$$

\subsubsection{Errore dovuto al troncamento}

\begin{align*}
    \frac{\left|x-\operatorname{fl}^T(x)\right|}{|x|} & \le\frac{\left|x-\operatorname{fl}^T(x)\right|}{\left|\operatorname{fl}^T(x)\right|}\le\frac{2^{e-b-n}}{2^{e-b}\cdot1.f} \\\\
                                                      & \le\frac{2^{e-b-n}}{2^{e-b}\cdot1.0}=\boxed{2^{-n}}
\end{align*}

Pertanto:
$$
    \boxed{
        \frac{\left|x-\operatorname{fl}^{T}(x)\right|}{\left|x\right|}\le2^{-n}}
$$
\subsubsection{Errore dovuto all'arrotondamento}
L'arrotondamento fa meglio perchè:
$$
    \left|x-\operatorname{fl}^{A}(x)\right|
    \le
    \frac12
    \left|\operatorname{succ}(\widetilde{x})-\widetilde{x}\right|=\frac{1}{2}\cdot2^{e-b-n}=2^{e-b-n-1}
$$

Dividendo per $|x|\ge2^{e-b}$ otteniamo:
$$
    \left|\frac{x-\operatorname{fl}^{A}(x)}{x}\right|\le\frac{2^{e-b-n-1}}{2^{e-b}}=2^{-n-1}
$$

Pertanto:
$$
    \boxed{\frac{\left|x-\operatorname{fl}^{A}(x)\right|}{\left|x\right|}\le2^{-n-1}}
$$

\section{Precisione di macchina}
\begin{defi}{Precisione macchina}{}
    Definiamo la \underline{precisione di macchina} come il più piccolo numero positivo $\epsilon$ tale che, se sommato a $1$ in virgola mobile, da un risultato diverso da $1$.
    $$
        \epsilon_{\mathrm{mach}}
        :=
        \min
        \left\{
        \epsilon>0:
        \operatorname{fl}(1+\epsilon)>1
        \right\}
    $$
\end{defi}

\subsection{Operazioni macchina}

Siano definite le operazioni esatte:

\begin{align*}
    \Delta:\  & \mathbb{R}\times\mathbb{R}\longrightarrow\mathbb{R},                 & (+,\cdot) \\
              & \mathbb{R}\times(\mathbb{R}\setminus\{0\})\longrightarrow\mathbb{R}, & (:)
\end{align*}

Operazione macchina associata \textcircled{$\Delta$}:
\begin{align*}
    \textcircled{$\Delta$}\quad & \mathbb{R}\times\mathbb{R}\longrightarrow\mathbb{F}                                                                                                      \\
                                & \left(x,y\right)\longrightarrow\textcircled{$\Delta$}\left(x,y\right)=\operatorname{fl}\left(\operatorname{fl}(x)\Delta\cdot \operatorname{fl}(y)\right)
\end{align*}

\subsection{Errore relativo della moltiplicazione}

$$
    \boxed{
        \operatorname{Err}_{\mathrm{rel}}(x\odot y)
        =
        \frac{
            |x\odot y-xy|
        }{
            |xy|
        }
        =
        \frac{
            \left|
            \operatorname{fl}
            \left(
            \operatorname{fl}(x)\operatorname{fl}(y)
            \right)
            -
            xy
            \right|
        }{
            |xy|
        }}
$$

\colorbox{yellow}{N.B.}: noi supponiamo che $$
    \frac{\left|x\cdot y-\operatorname{fl}(x\cdot y)\right|}{\left|x\cdot y\right|}\le\epsilon_{\text{mach}}
    \begin{cases}
         & 2^{-n}   \\
         & 2^{-n-1}
    \end{cases}
$$
$z=x\cdot y$

\begin{align*}
    \operatorname{Err}_{\mathrm{rel}}(x\odot y) & =\frac{\left|\operatorname{fl}\left(\operatorname{fl}(x)\cdot\operatorname{fl}(y)\right)-\operatorname{fl}(x\cdot y)+\operatorname{fl}(x\cdot y)-x\cdot y\right|}{\left|x\cdot y\right|}                                                                                                                                                                 \\
                                                & \le\frac{\left|\operatorname{fl}(\operatorname{fl}(x)\cdot\operatorname{fl}(y))-\operatorname{fl}(x\cdot y)\right|}{\left|x\cdot y\right|}+\overbrace{\frac{\left|\operatorname{fl}(x\cdot y)-x\cdot y\right|}{\left|x\cdot y\right|}}^{\epsilon_{\mathrm{mach}}}                                                                                        \\
                                                & =\frac{\overbrace{\left|\operatorname{fl}\left(\operatorname{fl}(x)\cdot\operatorname{fl}(y)\right)-\operatorname{fl}(x)\cdot\operatorname{fl}(y)\right|}^{\textcolor{red}{1}}+\overbrace{\left|\operatorname{fl}(x)\cdot\operatorname{fl}(y)-\operatorname{fl}(x\cdot y)\right|}^{\textcolor{blue}{2}}}{\left|x\cdot y\right|}+\epsilon_{\mathrm{mach}} \\
                                                & \le^{\textcolor{red}{1}}\frac{\left|\operatorname{fl}\left(\operatorname{fl}(x)\cdot\operatorname{fl}(y)\right)-\operatorname{fl}(x)\cdot\operatorname{fl}(y)\right|}{\left|\operatorname{fl}(x)\cdot\operatorname{fl}(y)\right|}\cdot\frac{\left|\operatorname{fl}(x)\cdot\operatorname{fl}(y)\right|}{\left|x\cdot y\right|}                           \\
                                                & \le \epsilon_{\mathrm{mach}}\cdot \frac{\left|\operatorname{fl}(x)\cdot\operatorname{fl}(y)\right|}{\left|x\cdot y\right|}                                                                                                                                                                                                                               \\
                                                & \le^{\textcolor{blue}{2}}\frac{\left|\operatorname{fl}(x)\cdot\operatorname{fl}(y)-x\cdot y+x\cdot y-\operatorname{fl}(x\cdot y)\right|}{\left|x\cdot y\right|}                                                                                                                                                                                          \\
                                                & \le\frac{\left|\operatorname{fl}(x)\cdot\operatorname{fl}(y)-x\cdot y\right|}{\left|x\cdot y\right|}+\frac{\left|\operatorname{fl}(x\cdot y)-x\cdot y\right|}{\left|x\cdot y\right|}\le \epsilon_{\mathrm{mach}}                                                                                                                                         \\
                                                & \le\epsilon_{\mathrm{mach}}+\frac{\left|\operatorname{fl}(x)\cdot\operatorname{fl}(y)-x\cdot y\right|}{\left|x\cdot y\right|}                                                                                                                                                                                                                            \\
    \textcolor{red}{1}+\textcolor{blue}{2}      & \le\epsilon_{\mathrm{mach}}\left(1+\frac{\left|\operatorname{fl}(x)\cdot\operatorname{fl}(y)\right|}{\left|x\cdot y\right|}\right)+\frac{\left|\operatorname{fl}(x)\cdot\operatorname{fl}(y)-x\cdot y\right|}{\left|x\cdot y\right|}                                                                                                                     \\
                                                & \le\epsilon_{\mathrm{mach}}\left(1+\epsilon_{\mathrm{mach}}(\underbrace{\dots\dots\dots}_{\text{dipende da $x,y$}})\right)
\end{align*}

Al primo ordine, il comportamento dell'errore relativo non dipende dai valori specifici di $x$ e $y$.\\
L'\underline{errore relativo della moltiplicazione macchina} $u$: $\epsilon+M\epsilon$.\\
La moltiplicazione è quindi considerata \colorbox{yellow}{stabile}, ossia l'errore relativo in output è controllato da quello in input.

\subsection{Errore relativo della somma}

L'errore relativo della somma macchina è:

\[
    \boxed{
        \operatorname{Err}_{\mathrm{rel}}(x\oplus y)
        =
        \frac{
            \left|
            \operatorname{fl}
            \left(
            \operatorname{fl}(x)+\operatorname{fl}(y)
            \right)
            -
            (x+y)
            \right|
        }{
            |x+y|
        }
        =
        \frac{
            |x\oplus y-(x+y)|
        }{
            |x+y|
        }}
\]

Se $x$ e $y$ hanno segno opposto e sono molto vicini tra loro, può verificarsi una forte \underline{cancellazione numerica}, ossia la perdita di cifre significative.

\begin{ex}{Cancellazione Numerica}{}
    Ipotesi in cui $x$ e $y$ si annullano fra loro lasciando solo un $\epsilon_{\mathrm{mach}}$:
    \[
        y=-x+\epsilon_{\mathrm{mach}}
    \]
\end{ex}
\colorbox{yellow}{Svolgimento}:
\begin{align*}
    \operatorname{Err}_{\mathrm{rel}}(x\oplus y) & \le\epsilon_{\mathrm{mach}}\left(1+\frac{|x|+|y|}{|x+y|}+\epsilon_{\mathrm{mach}}\frac{|x|+|y|}{|x+y|}\right)                        \\
                                                 & =\epsilon_{\mathrm{mach}}\left(1+\frac{2|x|+\epsilon_{\mathrm{mach}}}{\epsilon_{\mathrm{mach}}}+2|x|+\epsilon_{\mathrm{mach}}\right) \\
                                                 & =\left(\epsilon_{\mathrm{mach}}+2|x|+\epsilon_{\mathrm{mach}}+2|x|\epsilon_{\mathrm{mach}}+\epsilon_{\mathrm{mach}}^2\right)         \\
                                                 & =2\left(|x|+\epsilon_{\mathrm{mach}}\right)+\dots
\end{align*}
I $\dots$ assorbono $\epsilon_{\mathrm{mach}}^2$ e $2|x|\epsilon_{\mathrm{mach}}$ che sono numeri talmente piccoli rispetto agli altri da essere irrilevanti.\\
Significa che se $x$ è grande, anche l'errore relativo è enorme.\\

\subsection{Stabilità degli algoritmi}

\begin{ex}{Calcolo dell'integrale col metodo iterativo}{}
    \[I_n= \frac{1}{e}\int_{0}^{1}e^x\cdot x^n\cdot dx\]
\end{ex}
\colorbox{yellow}{Svolgimento}:
Conoscendo il valore di $I_0$ ossia: \[I_0=\frac{1}{e}\int_{0}^{1}e^x\cdot dx=\frac{1}{e}\left[e^x\right]^1_0=\frac{e-1}{e}\]

Inoltre si può integrare per parti $I_n$:
\begin{align*}
    I_n & =\frac{1}{e}\left[\left(e^xx^n\right)^1_0-\int^1_0 e^x\cdot nx^{n-1}\cdot dx\right] \\
        & =\frac{1}{e}\left[e-n\int^1_0 e^x\cdot x^{n-1}\cdot dx\right]                       \\
        & =1-n\frac{1}{e}\int_{0}^{1}e^x\cdot x^{n-1}\cdot dx                                 \\
        & =1-n\cdot I_{n-1}
\end{align*}
Di conseguenza conosco la \underline{successione definita per ricorsione}:
\[\boxed{\begin{cases}
             & I_0=\frac{e-1}{e}    \\
             & I_n=1-n\cdot I_{n-1}
        \end{cases}}\]
In realtà $I_{n-1}$ viene calcolato con un errore, quindi supponiamo di aver calcolato $\widetilde{I}_{n-1}=I_{n-1}+\epsilon_{n-1}$.
Quindi per calcolare $I_n$:
\begin{align*}
    I_n             & =1-n\cdot I_{n-1}                                                     \\
    \widetilde{I}_n & =1-n\cdot \widetilde{I}_{n-1}=\underbrace{1-nI_{n-1}}-n\epsilon_{n-1} \\
                    & =\boxed{I_n-n\epsilon_{n-1}}
\end{align*}
Questo risultato è iterabile dato che si verifica per ogni $n$, inoltre so che $\epsilon_n=-n\epsilon_{n-1}$.\\
Iterando $n$ volte ottengo:\[\boxed{\epsilon_n=(-1)^n\cdot n!\cdot\epsilon_0}\]

L'errore assoluto è quindi:\[|\widetilde{I}_n-I_n|=n!\cdot |\epsilon_0|=n!\cdot \operatorname{fl}\left(\frac{e-1}{e}\right)-\frac{e-1}{e}\]

Da cui ottengo:\[=n!\cdot\underbrace{\frac{|\operatorname{fl}\left(\frac{e-1}{e}\right)-\frac{e-1}{e}|}{\frac{e-1}{e}}}_{\le\epsilon_{\mathrm{mach}}}\cdot \underbrace{\frac{e-1}{e}}_{\le 1}\]
La frazione è l'errore relativo di rappresentazione di un numero reale\\
L'\underline{errore assoluto} è dominato da:
\[\boxed{\le n!\cdot\epsilon_{\mathrm{mach}}}\]

Adesso stimiamo quanto è $I_n$:
\[I_n=\frac{1}{e}\int_{0}^{1}x^ne^x\cdot dx=O(\frac{1}{n})\]
Dimostriamo che è vero:

Posso stimarlo da sopra con un integrale $\le$
\[nI_n=\frac{1}{e}\int_{0}^{1}nx^ne^x\cdot dx\le\frac{1}{e}\int_{0}^{1}enx^n\cdot dx=\frac{n}{n+1}\left[x^{n+1}\right]_0^1=\frac{n}{n+1}\]

E posso stimarlo da sotto con un'integrale $\ge$
\[nI_n=\frac{1}{e}\int_{0}^{1}nx^ne^x\cdot dx \ge \frac{1}{e}\int_{0}^{1}nx^n\cdot dx=\frac{n}{n+1}\cdot\frac{1}{e}\]
Abbiamo quindi scoperto che $I_n \approx \frac{1}{n+1}$\\
Adesso possiamo dire che l'\underline{errore relativo}:
\[\frac{|\widetilde{I}_n-I_n|}{|I_n|}\le \frac{n!\epsilon_{\mathrm{mach}}}{\frac{1}{n+1}}=(n+1)n!\epsilon_{\mathrm{mach}}=\boxed{(n+1)!\epsilon_{\mathrm{mach}}}\]\\
Questo algoritmo si dice \textbf{instabile}.\\

\begin{defi}{Algoritmo Stabile}{}
    Un algoritmo si dice \textbf{stabile} se non amplifica in modo incontrollato gli errori presenti sui
    parametri di input e i dati. In generale, si dice \textbf{stabile} un algoritmo per cui vale:

    \[
        \operatorname{Err}^{\mathrm{rel}}(\text{output})\le C_{\mathrm{stab}}\operatorname{Err}^{\mathrm{rel}}(\text{input})
    \]
    Dove $C_{\mathrm{stab}}$ è la \underline{costante di stabilità}.
\end{defi}

\subsection{Condizionamento dei problemi}

Per ogni problema abbiamo $x^*$ che è la soluzione, su $F:\ X\times D\to\mathbb{R}\ (\text{o }\mathbb{R}^m)$ dove, $X$ è il dominio della soluzione
e $D$ è l'insieme dei dati ammissibili.
\begin{defi}{Problema ben posto}{}
    Un problema si dice \textbf{ben posto} se:
    \begin{enumerate}
        \item \(\forall d\in D\quad\exists\ x^*\in X:\ \operatorname{F}(x^*,d)=0\) esiste una \textbf{soluzione}
        \item \(x^*\), la soluzione è \textbf{unica}
        \item $x^*$ varia con \textbf{continuità} rispetto ai dati\\
    \end{enumerate}
    La soluzione è data dalla \textbf{mappa soluzione}:\[d\mapsto x^*(d)\ :\ \operatorname{F}(x^*(d),d)=0\quad x^*\ :\ D\to X\]
\end{defi}
I \textbf{problemi inversi}, ossia quando dato l'output si vuole ricavare l'input, \textbf{non} sono ben posti; dato che più input possono avere lo stesso output.
Vengono quindi \textbf{regolarizzati}, ossia si modifica il problema in modo da penalizzare soluzioni "non fisiche" o "di bassa energia".

\begin{defi}{Condizionamento Numerico}{}
    È dato dal rapporto tra l'errore commesso sul risultato di un calcolo e l'incertezza sui dati in ingresso.\\
    Un problema si dice \textbf{ben condizionato} se la soluzione del problema con delle piccole variazioni, non cambia
    di molto la soluzione del problema originale. Un problema è \textbf{mal condizionato} se la soluzione è molto sensibile
    a piccole variazioni dei dati iniziali.
\end{defi}

\subsubsection{Condizionamento di un problema ben posto}
\begin{align*}
    \frac{\text{cond. locale}}{\text{assoluto}}\quad & \operatorname{k}(d):=\limsup_{\widetilde{d}\to d}\frac{\left|x(d)-x(\widetilde{d})\right|}{\left|d-\widetilde{d}\right|}
\end{align*}
Quindi $\operatorname{k}(d)$ è rappresentato dal limite da i dati approssimati $\widetilde{d}$ verso il dato vero $d$ del condizionamento assoluto diviso l'errore sul dato.
Questo limite rappresenta il rapporto incrementale, come nelle derivate.
\begin{align*}
    \limsup_{\widetilde{d}\to d}f(\widetilde{d}) & =\lim_{r\to 0^+}\sup_{\widetilde{d}\in\ ]d-r,d+r[}f(\widetilde{d})
\end{align*}
Ossia prendo l'intorno del dato $d$ e trovo il sup diminuendo $r$, in questo modo mi avvicino a $d$ tenendo $f$ più grande possibile.

$\operatorname{k}(d)$ è la pendenza del grafico di $x(d)$ \textbf{massima in modulo}. Se abbiamo una pendenza $\operatorname{k}(d)$, allora
vuol dire che  $\forall\ \epsilon>0\ \exists r>0$ tale che:$$\frac{\left|x(\widetilde{d})-x(d)\right|}{|\widetilde{d}-d|}<\operatorname{k}(d)+\epsilon\quad \forall\widetilde{d}\in\ ]d-r,d+r[$$
È quindi una stima dall'alto di $d$ che vale \underline{localmente}, ossia su un'intorno di $d$. Per questo si usa il $\limsup$.

Si verifica un problema quando $\operatorname{k}(d)>>1$ (quindi grande rispetto al contesto). Un \textbf{errore assoluto} piccolo su $d$ può
implicare un \textbf{errore assoluto} enorme su $x$. \textbf{Problema mal condizionato}.

\colorbox{yellow}{Grafico}:

\begin{tikzpicture}
    \begin{axis}[axis lines=middle, ymax=8, ymin=-8, xlabel=$x$, ylabel=$y$]
        \addplot[color=black, samples=100, domain=-2:1.5]{x^3+2};
        \addplot[color=black, samples=100, domain=1.5:3.5]{-(2*x-6)^2+4};
        \addplot[color=red, mark=o] coordinates {(1.5, 5.375)};
        \addplot[color=black, mark=*] coordinates {(1.5, -5)};
        \addplot[color=red, mark=*] coordinates {(3.5, 3)};
    \end{axis}
\end{tikzpicture}

\subsubsection{Condizionamento locale relativo}
$$\operatorname k^{\mathrm{rel}}(d):=\boxed{\limsup_{\widetilde{d}\to d}\underbrace{\frac{\left|x(\widetilde{d})-x(d)\right|}{\left|x(d)\right|}}_{\text{pendenza mappa sol.}}\cdot\frac{\left|d\right|}{\left|\widetilde{d}-d\right|}}$$

Ci sono problemi in cui il condizionamento non dipende dal dato $d$, ma è uguale dappertutto.

Può succedere che $\exists K\ (K^{\mathrm{rel}})$ tale per cui $\operatorname{k}(d)\le K\ \left(\operatorname{k}^{\mathrm{rel}}(d)\le K^{\mathrm{rel}}\right)\ \forall d\in D$. Dove $D$ è l'insieme di dati ammissibili.
In tal caso si parla di condizionamento assoluto/relativo \textbf{globale} e si ha:
\[\operatorname{k}=\sup_{d,\widetilde{d}\in D}\frac{\left|x(\widetilde{d})-x(d)\right|}{|\widetilde{d}-d|}\quad \blacktriangle \]

\begin{tikzpicture}
    \begin{axis}[axis lines=center, ymax=100, ymin=0, title={No cond. globale}, xlabel=$x$, ylabel=$y$]
        \addplot[color=black, samples=80, domain=0:6]{(1/(x-3))^2};
        \draw [black, dashed] (axis cs:3,\pgfkeysvalueof{/pgfplots/ymin}) -- (axis cs:3,\pgfkeysvalueof{/pgfplots/ymax});
    \end{axis}
\end{tikzpicture}\hskip 10pt
\begin{tikzpicture}
    \begin{axis}[title={Si cond. globale}, axis lines=center, domain = 0.01:6, samples = 200, ylabel=$y$, xlabel=$x$]
        \addplot[color=black, samples=80]{x^(3-x)};
        \addplot[color=red, domain=0:2, dashed]{2*x-1};
    \end{axis}
\end{tikzpicture}

$$\blacktriangle\quad \left|x(\widetilde{d})-x(d)\right|\le\operatorname{K}\left|\widetilde{d}-d\right|\quad \forall\ d,\widetilde{d}\in D$$
Ho un fattore moltiplicativo nel caso dell'errore assoluto.\\
Mentre per il relativo, l'errore varia per $\epsilon_{\mathrm{mach}}$ moltiplicato a $\operatorname{K}^{\mathrm{rel}}$:
\[\frac{\left|x(\widetilde{d})-x(d)\right|}{\left|x(d)\right|}\le\operatorname{K}^{\mathrm{rel}}\frac{\left|\widetilde{d}-d\right|}{|d|}\]

Se $\operatorname{K}^{\mathrm{rel}} \approx \frac{1}{\epsilon_{\mathrm{mach}}}\Rightarrow\frac{\left|x(\widetilde{d})-x(d)\right|}{\left|x(d)\right|}\approx 1=100\%\quad$
Quindi un'errore relativo di circa $1$ vuol dire che la soluzione è completamente diversa da quella vera.

\colorbox{yellow}{Esempio}:

\begin{tikzpicture}
    \begin{axis}[axis lines=center, domain=-2:12, xtick=\empty ,ytick=\empty]
        \addplot[color=black, samples=100]{-(1/2*x)+4};
        \addplot[color=green, samples=100]{-(1/2*x)+5};
        \node[anchor=south west, color=green] at (axis cs:10,0) {$x(\widetilde{q})$};
        \node[anchor=south west, color=black] at (axis cs:6.5,-1) {$x(q)$};
        \node[anchor=south east, color=black] at (axis cs:0,3.5) {$q$};
        \node[anchor=south west, color=green] at (axis cs:0,5) {$\widetilde q$};
        \draw [red] (axis cs:0,4) -- (axis cs:0,5);
        \draw [red] (axis cs:8,0) -- (axis cs:10,0);
    \end{axis}
\end{tikzpicture}

Con il grafico si nota come un minore errore su $q$ crea un errore maggiore su $x(q)$.

\subsubsection{Condizionamento globale di ordine $\alpha$}
\begin{ex}{Cosa succede quando la pendenza "esplode"?}{}
    Dato $d\in D$, la cui mappa soluzione è $d\mapsto x(d)$ con $D=[0,+\infty[$ e $x(d)=\sqrt{d}$.
\end{ex}

\begin{tikzpicture}
    \begin{axis}[axis lines=center, xmin=-2, ymin=-2, xlabel=$x$, ylabel=$y$, xtick=\empty, ytick=\empty]
        \addplot[samples=200]{x^(1/2)};
    \end{axis}
\end{tikzpicture}

In questo la caso la pendenza vicino a $0$ tende a $\infty$.\\
Quindi: \[\frac{\left|\sqrt{\widetilde{d}}-\sqrt{d}\right|}{\left|\widetilde{d}-d\right|}=\frac{\left|\sqrt{\widetilde{d}}-\sqrt{d}\right|}{\left|\left(\sqrt{\widetilde{d}}-\sqrt{d}\right)\left(\sqrt{\widetilde{d}}+\sqrt{d}\right)\right|}=\frac{1}{\left|\sqrt{\widetilde{d}}+\sqrt{d}\right|}\]
Ossia abbiamo che:
\[
    \limsup_{\widetilde{d}\to d}\frac{1}{2\sqrt{d}}
    \begin{cases}
         & d\ne0\quad\checkmark \\
         & d=0\quad+\infty
    \end{cases}
\]
Quindi avvicinandomi a $d=0$ non riesco a fare delle stime quantitative.\\
Si introduce quindi il \textbf{condizionamento globale assoluto} di \textbf{ordine} $\alpha\in(0,1)$
\[\operatorname{k}_{\alpha}=\boxed{\sup_{d,\widetilde{d}}\frac{\left|x(\widetilde{d})-x(d)\right|}{\left|\widetilde{d}-d\right|^\alpha}}\]
Lo usiamo per studiare problemi \textbf{molto mal condizionati}.

Nell'esempio $\alpha=\frac{1}{2}$:
\begin{align*}
     & =\frac{\left|\sqrt{\widetilde{d}}-\sqrt{d}\right|}{\left|\widetilde{d}-d\right|^{\frac 12}}                                                                                                          \\
     & =\frac{\epsilon}{\left(\epsilon(2\sqrt{d}+\epsilon)\right)^{\frac 12}}=\frac{\epsilon}{\left(\epsilon^2+2\epsilon\sqrt d\right)^{\frac 12}}\quad \text{se }d=0\rightarrow\frac{\epsilon}{\epsilon}=1
\end{align*}
se $d>0$ si mostra che è limitata.

Se vale:
\begin{align*}
     & \sup_{d,\widetilde{d}}\frac{\left|x(\widetilde{d})-x(d)\right|}{\left|\widetilde{d}-d\right|^\alpha}\le \operatorname{k}_\alpha  \\
     & \left|x(\widetilde{d})-x(d)\right|\le\operatorname{k}_\alpha\left|\widetilde{d}-d\right|^\alpha\quad\forall d,\widetilde{d}\in D
\end{align*}
\colorbox{yellow}{Nota}: $\alpha<1\quad\left|\widetilde{d}-d\right|<<1\Rightarrow\left|\widetilde{d}-d\right|^\alpha >>\left|\widetilde{d}-d\right|$.


\section{Algebra Lineare Numerica}

\subsection{Soluzione di Sistemi Quadrati}
\colorbox{yellow}{Problema}:\(\qquad A\in M_{n\times n}(\mathbb{R}),b\in\mathbb{R}^n \quad\) con $A$ invertibile.\\
Trovare $x\in\mathbb{R}^m$ tale che $\underbrace{Ax - b=0}_{\operatorname{F}(x,d)=0}$ funzione che prende un vettore reale e ne ritorna un'altro.\\
Dominio di $F\quad \underset{x}{\mathbb{R}^n}\times \underset{A}{M_{n\times n}(\mathbb{R})}\times \underset{b}{\mathbb{R}^n}\to \underset{Ax-b}{\mathbb{R}^n}$

\colorbox{yellow}{Osservazione}: $A$ invertibile $\Rightarrow$ esiste ed è unica $x:\ Ax=b$, quindi la mappa soluzione è:
\[x=A^{-1}b\]

Manca sapere se la mappa soluzione è continua, serve quindi una \underline{norma} sul dominio e una sul codominio.
Su $\mathbb{R}^n$ prendiamo una norma qualsiasi.

\subsection{Norma indotta}
\begin{defi}{Norma Indotta}{}
    Se $\|\cdot\|$ è una norma su $\mathbb{R}^m$ esiste una norma detta \textbf{naturale} su $M_{m\times m}(\mathbb{R})$ \textbf{indotta}
    da $\|\cdot\|$: \[\|A\|_*=\sup_{x\in\mathbb R^m\setminus\{0\}}\frac{\|Ax\|}{\|x\|}\]
    Questa norma indicata è $\|\cdot\|_*$
\end{defi}

\colorbox{yellow}{Proposizione}: $\|\cdot\|_*$ è una norma.

\colorbox{yellow}{Dimostrazione}: Devo dimostrare che rispetta le proprietà di una norma.
\begin{enumerate}
    \item Omogeneità di grado $1$. Prendiamo $\lambda\in\mathbb R\setminus\{0\}$ allora:
          \begin{align*}
              \|\lambda A\|_*\quad & =\sup_{x\in\mathbb R\setminus\{0\}}\frac{\|\lambda (Ax)\|_*}{\|x\|}=\sup_{x\in\mathbb R\setminus\{0\}}|\lambda|\cdot\frac{\|Ax\|}{\|x\|} \\
                                   & =|\lambda|\sup_{x\in\mathbb R\setminus\{0\}}\frac{\|Ax\|}{\|x\|}=|\lambda|\cdot\|A\|_*\quad \checkmark
          \end{align*}

    \item Una norma è $=0$ se e solo se è$\|0\|$. Quindi supponiamo che la norma sai uguale a $0$:
          \begin{align*}
              0=                                        & \|A\|_*=\sup_{x\in\mathbb R\setminus\{0\}}\frac{\|Ax\|}{\|x\|}\Rightarrow\|Ax\|=0\quad\forall\ x\in\mathbb R^m \\
                                                        & Ax=\begin{pmatrix}
                                                                 0      \\
                                                                 \vdots \\
                                                                 0
                                                             \end{pmatrix}\quad \forall\ x\in\mathbb R^m\qquad\text{Quindi è un vettore nullo}                           \\
              \underset{i\text{-esima colonna}}{A(:,i)} & =\quad A_{ei}=\begin{pmatrix}
                                                                            0      \\
                                                                            \vdots \\
                                                                            0
                                                                        \end{pmatrix}\quad\forall\ i\quad\Rightarrow\quad A=\begin{pmatrix}
                                                                                                                                0      & \dots & 0      \\
                                                                                                                                \vdots &       & \vdots \\
                                                                                                                                0      & \dots & 0
                                                                                                                            \end{pmatrix}\quad\checkmark
          \end{align*}

    \item Disuguaglianza triangolare:
\end{enumerate}
Ho dimostrato che è una norma dato che rispetta tutte e tre le proprietà.

Che norma si usa su $V\times W$ con $V$ e $W$ che hanno le norme $\|\cdot\|_v,\|\cdot\|_w$?
\begin{align*}
     & \|(v,w)\|_{v\times w}:=\|v\|_v+\|w\|_w\quad\text{Per esempio} \\
     & \text{La mappa soluzione: }x(A,b)=A^{-1}b \text{ è continua?}
\end{align*}
$\forall\epsilon>0\ \exists\delta>0$ tale che $$\left\|x\left(\widetilde{A},\widetilde{b}\right)-x\left(A,b\right)\right\|<\epsilon\quad\text{se}\quad\left\|\left(\widetilde{A},\widetilde{b}\right)-\left(A,b\right)\right\|\underset{\mathbb{M}_{m\times m}(\mathbb R)\times \mathbb R^m}{<\delta}$$
\colorbox{yellow}{Osservazione}: È composizione di $A\mapsto A^{-1}$ continua sulle matrici invertibili.\\
E $\left(A^{-1},b\right)\mapsto A^{-1}b$ continua.\\
Dunque la mappa soluzione è continua in $(A,b)$ se $A$ è invertibile.

Come è condizionato il problema?
\begin{itemize}
    \item Considero il caso in cui $A$ sia fissa ed "esatta".
          \begin{align*}
               & b\mapsto x(A,b)=A^{-1}b                                                                                                                               \\
               & A(x+\delta x)=b+\delta b\quad x+\delta x\text{ soluzione del sistema, con termine noto perturbato}                                                    \\
               & Ax=b                                                                                                                                                  \\
               & \text{Ci interessa quindi l'errore relativo: }\frac{\|\delta x\|}{\|x\|}\le \dots\cdot\frac{\|\delta b\|}{\|b\|}                                      \\
               & \underset{=b}{Ax}+A\delta x=b+\delta b                                                                                                                \\
               & \delta x=A^{-1}\delta b\Rightarrow \text{ Se due vettori sono uguali hanno anche norme uguali: }\left\|\delta x\right\|=\left\|A^{-1}\delta b\right\| \\
               & \Rightarrow\quad\left\|\delta x\right\|=\left\|A^{-1}\delta b\right\|\le\left\|A^{-1}\right\|_*\left\|\delta b\right\|\quad\circledast
          \end{align*}

          \begin{note}{}{}
              Dato che
              \[\frac{\left\|Ax\right\|}{\left\|x\right\|}\le\sup_{y}\frac{\left\|Ay\right\|}{\left\|y\right\|}=\left\|A\right\|_*\quad \forall x\in\mathbb{R}^n\]
              Allora:
              \[\boxed{\left\|Ax\right\|\underset{\forall x\in\mathbb{R}^n}{\le}\left\|A\right\|_*\cdot\left\|x\right\|}\]
          \end{note}
          Su $\circledast$ sto misurando in termini assoluti, devo invece misurare in termini \underline{relativi} dividendo per $\left\|x\right\|$:
          \begin{align*}
               & \frac{\left\|\delta x\right\|}{\left\|x\right\|}\le \left\|A^{-1}\right\|_*\left\|\delta b\right\|\cdot\frac{1}{\left\|x\right\|}\ \text{\textcircled{$\vartriangle$}}\quad\text{Mi serve poter stimare da sopra $\left\|x\right\|\ge\dots$ qualcosa} \\
               & \circledast\ \left\|x\right\|\ge\frac{\left\|b\right\|}{\left\|A\right\|_*}\Leftarrow
              \begin{cases}
                   & Ax=b\Rightarrow \left\|Ax\right\|=\left\|b\right\|                        \\
                   & \left\|Ax\right\|\le\left\|A\right\|_*\left\|x\right\|\quad \text{Nota 1}
              \end{cases}                                                                                                                                                                             \\
               & \textcircled{$\vartriangle$}+\circledast\quad\frac{\left\|\delta x\right\|}{\left\|x\right\|}\le\left\|A^{-1}\right\|_*\left\|A\right\|_*\cdot\frac{\left\|\delta b\right\|}{\left\|b\right\|}\ \forall b\in\mathbb R^n                               \\
               & \text{Dove $\left\|A^{-1}\right\|_*\left\|A\right\|_*$ è l'indice di condizionamento globale del problema semplificato.}
          \end{align*}

    \item Ora considero di perturbare anche $A$:
          \begin{align*}
               & \begin{cases}
                      & (A+\delta A)(x+\delta x)=b+\delta b \\
                      & Ax=b
                 \end{cases}                                     \\
               & (A+\delta A)\delta x=b+\delta b-(A+\delta A)x=\delta b-\delta Ax\quad\odot
          \end{align*}
          Noto che se \(\frac{\left\|\delta A\right\|_*}{\left\|A\right\|_*}<\frac{1}{\left\|A\right\|_*\left\|A^{-1}\right\|_*}\), allora
          \(\mathbb 1+A^{-1}\delta A\) è \textbf{invertibile}.

          \colorbox{yellow}{Dimostrazione}: Suppongo per assurdo che $\mathbb 1+A^{-1}\delta A$ non sia invertibile
          \begin{align*}
               & \exists y\ne 0\text{ tale che}\quad\left(\mathbb{1}+A^{-1}\delta A\right)y=0\quad\text{cioè}\quad A^{-1}\delta Ay=-y                                                                                                \\
               & \Rightarrow\ \left\|A^{-1}\delta Ay\right\|=\left\|y\right\|                                                                                                                                                        \\
               & 1=\frac{\left\|A^{-1}\delta Ay\right\|}{\left\|y\right\|}\le\sup_{x\in\mathbb R^n\ x\ne0}\frac{\left\|A^{-1}\delta Ax\right\|}{\left\|x\right\|}\underset{\text{Nota 1}}{=}\left\|A^{-1}\delta A\right\|_*          \\
               & 1\le\left\|A^{-1}\delta A\right\|_*\underset{\text{Nota 1}}{\le}\left\|A^{-1}\right\|_*\left\|\delta A\right\|_*=\left\|A^{-1}\right\|_*\left\|A\right\|_*\cdot\frac{\left\|\delta A\right\|_*}{\left\|A\right\|_*} \\
               & \frac{1}{\left\|A\right\|_*\left\|A^{-1}\right\|_*}\le\underset{\text{Perturbazione rel.}}{\frac{\left\|\delta A\right\|_*}{\left\|A\right\|_*}}\quad\text{\textbf{Assurdo} contraddice l'ipotesi.}\quad\checkmark
          \end{align*}

          Da $\odot$ segue che:
          \begin{align*}
                       & (A+\delta A)\delta x=\delta b-\delta Ax                                \\
                       & A(\mathbb 1+A^{-1}\delta A)\delta x=\delta b-\delta Ax                 \\
              \delta x & =(\mathbb 1+A^{-1}\delta A)^{-1}\cdot A^{-1}(\delta b-\delta Ax)       \\
                       & =(\mathbb 1+A^{-1}\delta A)^{-1}\cdot A^{-1}(\delta b-\delta AA^{-1}b)
          \end{align*}
          Passiamo alle norme
          \begin{align*}
              \left\|\delta x\right\| & =\dots\le\left\|\left(\mathbb 1+A^{-1}\delta A\right)^{-1}\right\|_*\left\|A^{-1}\right\|_*\left\|\delta b-\delta AA^{-1}b\right\|                                                         \\
                                      & \le\left\|\left(\mathbb 1+A^{-1}\delta A\right)^{-1}\right\|_*\left\|A^{-1}\right\|_*\left(\left\|\delta b\right\|+\left\|\delta AA^{-1}b\right\|\right)                                   \\
                                      & \le\left\|\left(\mathbb 1+A^{-1}\delta A\right)^{-1}\right\|_*\left\|A^{-1}\right\|_*\left(\left\|\delta b\right\|+\left\|\delta A\right\|_*\left\|A^{-1}\right\|_*\left\|b\right\|\right)
          \end{align*}

          Prima abbiamo usato $\frac{1}{\left\|x\right\|}\le\frac{\left\|A\right\|_*}{\left\|b\right\|}$. Quindi Per l'errore relativo:
          \begin{align*}
               & \frac{\left\|\delta x\right\|}{\left\|x\right\|}\le\left\|\left(\mathbb 1+A^{-1}\delta A\right)^{-1}\right\|_*\left\|A^{-1}\right\|_*\left\|A\right\|_*\left(\frac{\left\|\delta b\right\|}{\left\|b\right\|}+\left\|A^{-1}\right\|_*\left\|\delta A\right\|_*\right)                                                                              \\
               & \text{$\left\|\delta A\right\|_*$ è un errore assoluto, noi stiamo usando solo errori relativi, quindi divido e moltiplico per $\left\|A\right\|_*$}                                                                                                                                                                                               \\
               & \frac{\left\|\delta x\right\|}{\left\|x\right\|}\le\left\|\left(\mathbb 1+A^{-1}\delta A\right)^{-1}\right\|_*\boxed{\left\|A^{-1}\right\|_*\left\|A\right\|_*}\left(\frac{\left\|\delta b\right\|}{\left\|b\right\|}+\boxed{\left\|A^{-1}\right\|_*\left\|A\right\|_*}\frac{\left\|\delta A\right\|_*}{\left\|A\right\|_*}\right)\quad\circledast
          \end{align*}
          \begin{defi}{Numero di condizionamento}{}
              La quantità $\left\|A^{-1}\right\|_*\left\|A\right\|_*$ è detta \underline{numero di condizionamento} della matrice $A$ rispetto alla norma $\left\|\cdot\right\|$
              si denota con $k(A,\left\|\cdot\right\|)$ o $\operatorname{Cond}(A,\left\|\cdot\right\|)$
          \end{defi}

          La stima $\circledast$ è una stima di condizionamento \textbf{relativo locale}.\\
\end{itemize}

Tipicamente per matrici grandi $\begin{cases}
         & \text{Piene}\quad n>10^4  \\
         & \text{Sparse}\quad n>10^6 \\
    \end{cases}$\\
$\operatorname{Cond}(A,\left\|\cdot\right\|)$ è facilmente $\approx 10^4\ 10^5$ in casi buoni.\\
Sbagliare $A$ di $10^{-16}$ relativo, implica avere un errore relativo su $x$ di $10^{-8}$. Quindi su 16 cifre significative
ne calcolo solo 8 correttamente: questo è un caso buono per le matrici grandi.\\
\colorbox{yellow}{Salvezza}: In generale tutti i $\le$ mostrati finora sono in realtà stretti, per ottenere l'uguaglianza, ossia l'errore massimo
bisogna prendere il peggior dato con la peggior perturbazione (\textbf{molto improbabile}).

Riprendiamo il caso in cui $A$ è fissa ed esatta dove abbiamo $Ax=b$ ossia, perturbando $A$:
\begin{align*}
     & A\left(x+\delta x\right)=b+\delta b\ \to\ \frac{\left\|\delta x\right\|}{\left\|x\right\|}\le\operatorname{k}(A,\left\|\cdot\right\|)\frac{\left\|\delta b\right\|}{\left\|b\right\|} \\
     & \operatorname{k}(A,\left\|\cdot\right\|)=\left\|A\right\|_*\left\|A^{-1}\right\|_*
\end{align*}
Si tratta di una disequazione \underline{sharp}?  (potrebbe esistere il caso di uguaglianza)

Suppongo di usare $\left\|\cdot\right\|=\left\|\cdot\right\|_2$, ossia la \textbf{norma euclidea} $\left\|v\right\|_2=\left(\sum v_i^2\right)^{1/2}$.\\
Chi è la \textbf{Norma indotta} su $M_{n\times n}(\mathbb R)$?
\begin{align*}
    \left\|A\right\|_* & =\max_{x\ne 0}\frac{\left\|Ax\right\|_2}{\left\|x\right\|_2}=\left(\max_{x\ne 0}\frac{\sum\left(Ax\right)^2_i}{\sum x^2_i}\right)^{\frac 12} \\
                       & \text{Ricordiamo la proprietà del prodotto riga x colonna tra vettori: }\boxed{\sum y^2_i=y^T\cdot y}                                        \\
                       & =\left[\max_{x\ne0}\frac{x^TA^TAx}{x^Tx}\right]^{\frac 12}\quad\circledast\quad G=A^TA=G^T\text{ matrice simmetrica}
\end{align*}
\begin{note}{}{}
    $M$ è una matrice quadrata, \underline{se} è diagonalizzabile allora:$$M=P\Lambda P^{-1}$$
    Dove $\lambda$ è la matrice diagonale con gli autovalori di $M$ e $P$ è una matrice invertibile le cui colonne sono autovettori di $M$.
\end{note}

\subsubsection{Teorema Spettrale}
\begin{theo}{Teorema Spettrale}{}
    Ho $G\in M_{n\times n}(\mathbb R)$ \textbf{simmetrica} allora: $\exists\ \lambda_1,\lambda_2,\dots,\lambda_n$
    e $Q\in M_{n\times n(\mathbb R)}$ matrice \textbf{ortogonale} ($Q^TQ=QQ^T=\mathbb 1_n$) tali che:
    $$\boxed{G=Q\Lambda Q^T}\quad \Lambda=\begin{pmatrix}
            \lambda_1 & \dots  & 0         \\
            \vdots    & \ddots & \vdots    \\
            0         & \dots  & \lambda_n
        \end{pmatrix}$$
\end{theo}

Quindi:
\begin{align*}
     & x^TA^TAx=x^TGx\ge0\text{ dato che è uguale a }\left\|Ax\right\|^2_2                     \\
     & \text{È }=0\text{ solo se }Ax=0\text{ se }A\text{ è di rango pieno ciò implica che }x=0 \\
     & \Rightarrow \lambda_i>0\text{ se }A\text{ è invertibile}
\end{align*}
$G$ si dice quindi, \textbf{simmetrica strettamente definita positiva}, ossia tutti i suoi autovalori sono maggiori stretti di $0$.

\begin{align*}
    \circledast\left\|A\right\|_*^2 & =\max_{x\ne0}\frac{x^TGx}{x^Tx}\underset{\text{Teo. spett.}}{=}\max_{x\ne0}\frac{x^TQ\Lambda Q^Tx}{x^Tx}                      \\
                                    & =\max_{x\ne0}\frac{x^TQ\Lambda Q^Tx}{x^T\underset{=QQ^T}{\mathbb 1}x}=\max_{x\ne0}\frac{x^TQ\Lambda Q^Tx}{x^TQQ^Tx}           \\
    \text{Dato che }Q               & \text{ è invertibile, allora posso fare un cambio di variabile: }y=Q^Tx\quad x=0\Leftrightarrow y=0                           \\
                                    & =\max_{y\ne0}\frac{y^T\Lambda y}{y^Ty}=\max_{y\ne0}\left(\frac{y}{\left\|y\right\|}\right)^T\Lambda\frac{y}{\left\|y\right\|}
\end{align*}
Sono quindi passato da una generica matrice positiva ad un matrice \underline{diagonale} con diagonale positiva.
\begin{align*}
     & =max_{z:\ \left\|z\right\|_2=1}z^T\Lambda z                                                                                                                                                                                      \\
     & \text{Vogliamo, quindi, risolvere: }                                                                                                                                                                                             \\
     & max\sum_{i=1}^n \lambda_i z^2_i\quad\text{con}\quad\sum_{i=1}^n z_i^2=1                                                                                                                                                          \\
     & \text{Ossia voglio trovare il massimo autovalore di }G                                                                                                                                                                           \\
     & \text{Provo con: }z=e_{\underline{i}}\ :\ \lambda_{\underline{i}}\ge\lambda_i\ \forall i=1\dots n\quad\text{suppongo: }\lambda_1\ge\lambda_2\ge\dots\ge\lambda_n>0                                                               \\
     & \sum_{i=1}^{n}\lambda_iz^2_i=\underset{\text{Separo il primo termine}}{\lambda_1\left(z_1^2+\sum_{i=2}^{n}\frac{\lambda_i}{\lambda_1}z^2_i\right)}\le\lambda_1\underbrace{\left(z_1^2+\sum_{i=2}^{n}z^2_i\right)}_{=1}=\lambda_1 \\
     & \text{Se prendo }z=e_1\text{ (primo vettore della base canonica) ottengo:}                                                                                                                                                       \\
     & \sum_{i=1}^{n}\lambda_iz^2_i=\lambda_1\cdot1^2+\lambda_2\cdot0^2+\dots+\lambda_n\cdot0^2=\boxed{\lambda_1}
\end{align*}
Quindi vale:$$\left\|A\right\|_*^2=\max_{z:\ \left\|z\right\|_2=1}\frac{z^TGz}{z^Tz}=\lambda_{\mathrm{max}}(G)\text{ autovalore massimo}$$
$$\boxed{\left\|A\right\|_*=\sqrt{\lambda_{\mathrm{max}}\left(A^TA\right)}\text{ se }\left\|\cdot\right\|=\left\|\cdot\right\|_2}$$

\subsubsection{Decomposizione ai valori singolari (SVD)}
\begin{theo}{}{}
    Se $A\in M_{n\times n}(\mathbb R)$ di rango pieno allora $\exists\ U,\Sigma\ V$ con $U,V$ \textbf{ortogonali} $n\times n\ \Sigma=\begin{pmatrix}
            \sigma_1 & \dots  & 0        \\
            \vdots   & \ddots & \vdots   \\
            0        & \dots  & \sigma_n
        \end{pmatrix}
    $ \textbf{diagonale}
    con $\sigma_i>0$. Tale per cui:$$\boxed{A=U\Sigma V^T}$$
\end{theo}
Adesso vediamo cosa sono $U,\Sigma,V$:
\begin{align*}
    A^TA=V\Sigma U^T\cdot U\Sigma V^T=V\Sigma\cdot\Sigma V^T=V\Sigma^2 V^T\Rightarrow \sigma_i^2\text{ autovalore di }A^TA
\end{align*}
Si nota che le colonne di $V$ sono gli autovettori di $A^TA$. Per dimostrarlo, prendiamo una colonna di $V$, la moltiplichiamo e
vediamo se otteniamo $\lambda$ volte la stessa colonna (proprietà \textbf{autovalore}).\\
$V(:,j)$ è la $j$-esima colonna.
\begin{align*}
     & A^TA\cdot V(:,j)=V\Sigma^2\underbrace{V^TV(:,j)}_{w}                                                                                         \\
     & w_i=V^T(i,:)V(:j)=V(:,i)^TV(:,j)=\left(V^TV\right)(i,j)=\mathbb 1(i,j)\begin{cases}
                                                                                  & 1\quad\text{se }i=j    \\
                                                                                  & 0\quad\text{se }i\ne j
                                                                             \end{cases}                                              \\
     & A^TA\cdot V(:,j)=V\Sigma^2 e_j=V\sigma_j^2e_j=\sigma^2_jVe_j=\underset{\text{autovalore}}{\sigma^2_j}\ \underset{\text{autovettore}}{V(:,j)}
\end{align*}
Questo spiega cos'è $V$ e $\Sigma^2$. Ora vediamo cos'è la matrice $U$:
\begin{align*}
    AA^T=         & \ U\Sigma V^T\cdot V\Sigma U^T=U\Sigma^2 U^T                                           \\
    \text{Adesso} & \text{ non serve rifare i conti, sono gli stessi ma con nomi diversi. Quindi ottengo:} \\
    \Rightarrow   & \ \sigma_j^2\text{ sono anche autovalori di }AA^T                                      \\
    \Rightarrow   & \ \text{Gli autovettori di }AA^T\text{ sono le colonne di }U
\end{align*}

Adesso continuiamo con il condizionamento in norma $2$ e nel vedere se la disequazione è \underline{sharp}:
\begin{align*}
     & \left\|A\right\|_*=\sqrt{\lambda_{\mathrm{max}}\left(A^TA\right)}=\sqrt{\lambda_{\mathrm{max}}\left(V\Sigma^2V^T\right)}=\sqrt{\max \sigma_i^2}=\sigma_{\mathrm{max}} \\
     & \text{Per il condizionamento serve anche:}                                                                                                                            \\
     & \left\|A^{-1}\right\|_*=\sqrt{\lambda_{\mathrm{max}}\left(A^{-T}A^{-1}\right)}                                                                                        \\
     & A=U\Sigma V^T\Rightarrow A^{-1}=V\Sigma^{-1}U^T\qquad\Sigma^{-1}=\begin{pmatrix}
                                                                            \frac{1}{\sigma_1} & \dots  & 0                  \\
                                                                            \vdots             & \ddots & \vdots             \\
                                                                            0                  & \dots  & \frac{1}{\sigma_n}
                                                                        \end{pmatrix}                                                     \\
     & A^{-T}A^{-1}=U\Sigma^{-1}V^T\cdot V\Sigma^{-1}U^T=U\Sigma^{-2}U^T                                                                                                     \\
     & \left\|A^{-1}\right\|_*=\sqrt{\lambda_{\mathrm{max}}\left(U\Sigma^{-2}U^T\right)}=\max\frac{1}{\sigma_i}=\frac{1}{\sigma_{\mathrm{min}}}                              \\
     & \boxed{\operatorname{k}\left(A,\left\|\cdot\right\|_2\right)=\left\|A\right\|_*\left\|A^{-1}\right\|_*=\frac{\sigma_{\mathrm{max}}}{\sigma_{\mathrm{min}}}}
\end{align*}
$\left\|A\right\|_*$ nel caso $\left\|\cdot\right\|=\left\|\cdot\right\|_2$ su $\mathbb R^n$ viene denotata con $\left\|A\right\|_2$.\\

Torniamo a domandarci se:$$\frac{\left\|\delta x\right\|_2}{\left\|x\right\|_2}\le\operatorname{k}\left(A,\left\|\cdot\right\|_2\right)\frac{\left\|\delta b\right\|_2}{\left\|b\right\|_2}$$
è irrealistica, cioè esiste l'uguaglianza?\\
Per dimostrarla abbiamo usato 2 stime:
\begin{enumerate}
    \item $\left\|b\right\|_2=\left\|Ax\right\|_2\le\left\|A\right\|_2\left\|x\right\|_2\quad$Può valere nel caso di uguaglianza?
          \begin{align*}
               & \left\|Ax\right\|_2^2=x^TV\Sigma U^T\cdot U\Sigma V^Tx=x^TV\Sigma^2V^Tx \\
               & y=V^Tx\quad\Rightarrow\quad y^T\Sigma^2y                                \\
               & \left\|x\right\|_2^2=\left\|V^Tx\right\|_2^2=\left\|y\right\|_2^2
          \end{align*}
          \begin{note}{Ragionamento}{}
              Il $\max\frac{\left\|Ax\right\|_2^2}{\left\|x\right\|_2^2}$ è raggiunto se $x$ è proporzionale all'autovettore relativo all'autovalore massimo di $A^TA$, ossia $\sigma_{\mathrm{max}}^2$ cioè la prima colonna di $V$.
              Quindi esiste l'uguaglianza.
          \end{note}

          Se $x=c\cdot V(:,1)\quad Ax=b\Rightarrow b=c\cdot A^{-1}V(:,1)=c\cdot V\Sigma^{-1} U^TV(:,1)$ allora:
          \[\left\|b\right\|_2=\left\|A\right\|_2\left\|x\right\|_2\qquad\text{(esiste l'uguaglianza)}\]

    \item $\left\|\delta x\right\|_2=\left\|A^{-1}\delta b\right\|_2\le\left\|A^{-1}\right\|_2\left\|\delta b\right\|_2\quad$ Esiste l'uguaglianza?

          Usando nuovamente la \underline{nota 3} il caso di uguaglianza, si verifica se il vettore $\delta b$ è proporzionale all'autovettore relativo all'autovalore
          massimo di $A^{-T}A^{-1}\Rightarrow\frac{1}{\sigma_{\mathrm{min}}}$. Dunque $\delta b=\epsilon U(:,n)$.
\end{enumerate}

Quindi la stima è \underline{sharp}, cioè ottimale dato che esiste l'uguale. In generale, è una stima pessimistica.

\subsection{Risoluzione pratica $Ax=b$}
\colorbox{yellow}{Osservazione}:\\Se la matrice $A$ è triangolare superiore $\begin{pmatrix}
        1 & 1 & 1 \\
        0 & 1 & 1 \\
        0 & 0 & 1
    \end{pmatrix}$ o inferiore $\begin{pmatrix}
        1 & 0 & 0 \\
        1 & 1 & 0 \\
        1 & 1 & 1
    \end{pmatrix}$: So risolvere il problema.

\subsubsection{Caso inferiore}
\begin{align*}
     & a_{11}x_1=b_1\Rightarrow x_1=\frac{b_1}{a_{11}}\qquad\text{(Sto ignorando i valori $0$)}                                                                   \\
     & a_{21}x_1+a_{22}x_2=b_2\Rightarrow x_2=\frac{b_2-a_{21}x_1}{a_{22}}                                                                                        \\
     & \text{In generale: }x_1=\frac{b_1}{a_{1}}\quad x_k=\boxed{\frac{b_k-\sum_{j=1}^{k-1}a_{k,j}x_j}{a_{k,k}}}\quad\text{(algoritmo di sostituzione in avanti)} \\
\end{align*}

\subsubsection{Caso superiore}
\begin{align*}
     & a_{nn}x_n=b_n\Rightarrow x_n=\frac{b_n}{a_{nn}}                                                                                   \\
     & \text{In generale: }x_{k}=\boxed{\frac{b_k-\sum_{j=k+1}^{n}a_{k,j}x_j}{a_{k,k}}}\quad\text{(algoritmo sostituzione all'indietro)}
\end{align*}

\subsection{Fattorizzazione}
Dato che le matrici triangolari sono poche, si usa il metodo di \underline{eliminazione gaussiana}, per ottenerne di più.
\begin{defi}{Fattorizzazione}{}
    Il metodo di eliminazione gaussiana si può vedere come \textbf{fattorizazzione} di matrice. La fattorizzazione $LU$ tale che
    $A=LU$ dove $L$ sta per "lower triangular" invertibile e $U$ "upper triangular" invertibile.
\end{defi}
A cosa serve una fattorizzazione di questo tipo?

$$Ax=b\Leftrightarrow \underset{\text{associativo}}{LUx}=b$$
Faccio un cambio di variabile $y=Ux$ e ottengo:
\begin{itemize}
    \item $Ly=b$ posso calcolare $y$ con la sostituzione in avanti.
    \item $Ux=y$ posso calcolare $x$ con la sostituzione all'indietro.
\end{itemize}

Rendiamo algoritmico il \underline{metodo di eliminazione gaussiana} (MEG):
\begin{align*}
    \widetilde{\mathcal{L}}^{(k)}\left(m^{(k)}\right)=\mathbb 1-\begin{pmatrix}
                                                                    0 \\
                                                                    m^k
                                                                \end{pmatrix}e^T_k\quad m^{(k)}\in \mathbb R^{n-k}
\end{align*}
\begin{note}{}{}
    $\begin{pmatrix}
            0 \\
            m^k
        \end{pmatrix}$ è una matrice con tutti $0$ tranne la parte sotto la diagonale della $k$-esima colonna che ha da $m^k_1$ a $m^k_{n-k}$.\\
    Se moltiplicata con $e^T_k$, si ottiene la matrice:
    \[\begin{pmatrix}
            1         & 0         & \dots  & \dots     & 0      \\
            m_1^1     & 1         & 0      &           & \vdots \\
            \vdots    & m^2_1     & \ddots & \ddots    & \vdots \\
            \vdots    & \vdots    & \ddots & 1         & 0      \\
            m^1_{n-1} & m^2_{n-2} & \dots  & m_1^{n-1} & 1
        \end{pmatrix}\]
    È \textbf{triangolare superiore} e la chiamo $L$.
\end{note}
\[\widetilde{\mathcal{L}}^{(k)}\left(m^{(k)}\right)=\begin{pmatrix}
        1      & \dots    & 0      \\
        \vdots & 1        & \vdots \\
        0      & -m^{(k)} & 1
    \end{pmatrix}\]
Nella $k$-esima colonna ci metto $-m^{(k)}$ per questo è lunga $n-k$, sono gli elementi sotto la diagonale.

Cosa succede se faccio $\widetilde{\mathcal{L}}^{(k)}\left(m^{(k)}\right)\cdot A$?
\begin{align*}
    \widetilde{\mathcal{L}}^{(k)}\left(m^{(k)}\right)A & =\left(\mathbb 1_n-\begin{pmatrix}
                                                                                    0 \\
                                                                                    m^k
                                                                                \end{pmatrix}e^T_k\right)A\underset{prop. assoc.}{=}A-\begin{pmatrix}
                                                                                                                                      0 \\
                                                                                                                                      m^k
                                                                                                                                  \end{pmatrix}\underbrace{e^T_kA}_{k\text{-esima riga di }A} \\
                                                       & =A-\begin{pmatrix}
                                                                0 \\
                                                                m^k
                                                            \end{pmatrix}A(k,:)=A-\begin{pmatrix}
                                                                                      0               & \dots & 0               \\
                                                                                      0               & \dots & k\text{ righe}  \\
                                                                                      m^k_1A(k,1)     & \dots & m^k_1A(k,n)     \\
                                                                                      \vdots          &       & \vdots          \\
                                                                                      m^k_{n-k}A(k,1) & \dots & m^k_{n-k}A(k,n)
                                                                                  \end{pmatrix}
\end{align*}
Ossia ad ogni riga di $A$ da $k+1$ a $n$ togliamo la $k$-esima riga di $A$ moltiplicata per $m_1^k,\dots,m^k_{n-k}$ (\textbf{moltiplicatori}).
Nel MEG scelgo i moltiplicatori per eliminare tutta la parte sottodiagonale di $A$ nella colonna $k$-esima.

La colonna $k$-esima di $\widetilde{\mathcal{L}}^{(k)}\left(m^{(k)}\right)A$ è:
\[A(:,k)-\begin{pmatrix}
        0              \\
        k\text{ volte} \\
        0              \\
        m_1A(k,k)      \\
        \vdots         \\
        m_{n-k}^kA(k,k)
    \end{pmatrix}=\begin{pmatrix}
        A(1,k)                     \\
        \vdots                     \\
        A(k,k)                     \\
        \noindent\rule{3.5cm}{1pt} \\
        A(k+1,k)-m^k_1A(k,k)       \\
        \vdots                     \\
        A(n,k)-m^k_{n-k}A(k,k)
    \end{pmatrix}\]
La parte bassa è quella che voglio annullare. Per far ciò devo "scegliere" $m$ nel seguente modo:
\[\circledast\quad m^k_l=\frac{A(k+l,k)}{A(k,k)}\]
Devo verificare che $A(k,k)\ne0$, ossia la diagonale diversa da zero, perchè potrebbe esserci il caso:
\begin{ex}{Problema Diagonale $=0$}{}
    $A=\begin{pmatrix}
            0 & 1 \\
            1 & 0
        \end{pmatrix}\quad \operatorname{det}(A)=-1$ La matrice è invertibile nonostante la diagonale $=0$.
\end{ex}

\colorbox{yellow}{Ipotesi}:
Per il momento supponiamo che $A(k,k)\ne0$.\\
Poniamo $\widetilde{\operatorname{L}}^{(k)}=\widetilde{\mathcal{L}}^{(k)}\left(m^{(k)}\right)$ con $m^k$
scelto come in $\circledast$.

\subsubsection{Algoritmo Instabile}
\begin{itemize}
    \item[] $A^{(1)}=A$
    \item[] $A^{(k)}=\widetilde{\operatorname{L}}^{(k-1)}A^{(k-1)}$ Quello $k$-esimo è dato da quello $(k-1)$-esimo
\end{itemize}
Se vale l'ipotesi ad ogni passo: $\boxed{A^{(k-1)}(k-1,k-1)\ne0}$. Dopo $n-1$ passi avrò:
\[A^n=\widetilde{\operatorname{L}}^{(n-1)}A^{(n-1)}=\widetilde{\operatorname{L}}^{(n-1)}\widetilde{\operatorname{L}}^{(n-2)}A^{(n-2)}\dots=\widetilde{\operatorname{L}}^{(n-1)}\dots\widetilde{\operatorname{L}}^{(1)}A\]
La matrice $A^n$ sarà, quindi, una triangolare superiore; perchè al primo passo abbiamo preso la colonna $1$ di $A$ ed eliminato
la sua parte sotto la diagonale. Quindi $A^{(2)}$ ha la colonna $1$ con la parte sottodiagonale \textbf{nulla}.
Poi $A^{(3)}=\widetilde{\operatorname{L}}^{(2)}A^2=\left(\mathbb 1-\begin{pmatrix}
            0 \\
            m^2
        \end{pmatrix}e_2^T\right)A^2=A^2-\begin{pmatrix}
        0 \\
        m^2
    \end{pmatrix}e^T_2A^{(2)}$ e dato che $A^2$ ha la prima colonna sottodiagonale $0$ e viene sottratta ad una cosa che ha la prima colonna tutto $0$.
Quindi anche $A^{(3)}$ ha la prima colonna sottodiagonale $0$ e in più anche la seconda colonna ha la stessa proprietà per $\circledast$.\\
Iterando questo ragionamento otteniamo che $A^{(n)}$ ha le colonne $1,2,\dots,n-1$ che rispettano la stessa proprietà. Di conseguenza è \textbf{triangolare superiore}.

Pongo $U=A^{(n)}$, ossia $U=\widetilde{\operatorname{L}}^{(n-1)}\widetilde{\operatorname{L}}^{(n-2)}\dots\ \widetilde{\operatorname{L}}^{(1)}A\quad$ Dato che ogni $\widetilde{\operatorname{L}}^{(k)}$ è invertibile e voglio ottenere $LU=A$, divido per $L$:
\[\left(\widetilde{\operatorname{L}}^1\right)^{-1}\left(\widetilde{\operatorname{L}}^2\right)^{-1}\dots\left(\widetilde{\operatorname{L}}^{(n-1)}\right)^{-1}U=A\]

Adesso devo:
\begin{itemize}
    \item calcolare $\left(\widetilde{\operatorname{L}}^{(k)}\right)^{-1}$
    \item calcolare il prodotto
    \item mostrare che è triangolare inferiore
\end{itemize}

\begin{align*}
     & \widetilde{\operatorname{L}}^{(k)}=\widetilde{\mathcal{L}}^{(k)}\left(m^{(k)}\right) \text{ con }m^{(k)}\text{ scelto secondo } \circledast                                                              \\
     & \widetilde{\mathcal{L}}^{(k)}\left(m^{(k)}\right)\cdot\widetilde{\mathcal{L}}^{(k)}\left(-m^{(k)}\right)=\left(\mathbb 1-\begin{pmatrix}
                                                                                                                                        0 \\
                                                                                                                                        m^k
                                                                                                                                    \end{pmatrix}e^T_k\right)
    \left(\mathbb 1+\begin{pmatrix}
                            0 \\
                            m^k
                        \end{pmatrix}e^T_k\right)                                                                                                                                                                   \\
     & =\mathbb 1 + \begin{pmatrix}0\\m^k\end{pmatrix}e^T_k-\begin{pmatrix}0\\m^k\end{pmatrix}e^T_k-\underbrace{\begin{pmatrix}0\\m^k\end{pmatrix}e^T_k\begin{pmatrix}0\\m^k\end{pmatrix}e^T_k}_{\triangle}     \\
     & \begin{pmatrix}0\\m^k\end{pmatrix}e^T_k\begin{pmatrix}0\\m^k\end{pmatrix}e^T_k=\begin{pmatrix}0\\m^k\end{pmatrix}\underbrace{\left[e^T_k\begin{pmatrix}0\\m^k\end{pmatrix}\right]}_{\text{scalare}}e^T_k \\
     & \triangle=\left[e^T_k\begin{pmatrix}0\\m^k\end{pmatrix}\right]\begin{pmatrix}0\\m^k\end{pmatrix}e^T_k                                                                                                    \\
\end{align*}
\begin{align*}
     & \text{Cosa significa fare }e^T_k\begin{pmatrix}0\\m^k\end{pmatrix}\text{, vuol dire: }\begin{pmatrix}0&\dots&0&\underbrace{1}_k&0&\dots&0\end{pmatrix}
    \cdot\begin{pmatrix}
             0                  \\
             \vdots             \\
             \underbrace{0}_{k} \\
             m_1^k              \\
             \vdots             \\
             m_{n-k}^k
         \end{pmatrix}=0                                                                                                                                   \\
     & \triangle=0s \Rightarrow\ \widetilde{\mathcal{L}}^{(k)}\left(m^{(k)}\right)\cdot \widetilde{\mathcal{L}}^{(k)}\left(-m^{(k)}\right)=\mathbb 1
\end{align*}
Dunque, abbiamo dimostrato che una è l'inversa dell'altra, e
\[\operatorname{L}^{(k)}:=\left(\widetilde{\operatorname{L}}^{(k)}\right)^{-1}=\left(\widetilde{\mathcal{L}}^{(k)}\left(m^{(k)}\right)\right)=\left(\widetilde{\mathcal{L}}^{(k)}\left(-m^{(k)}\right)\right)=\mathbb 1+\begin{pmatrix}
        0 \\
        m^k
    \end{pmatrix}e^T_k\]

Adesso passo al secondo punto, ossia calcolare i prodotti:
\begin{align*}
     & \operatorname{L}^{k}\cdot\operatorname{L}^{k+1}=\left(\mathbb 1+\begin{pmatrix}
                                                                               0 \\
                                                                               m^k
                                                                           \end{pmatrix}e^T_k\right)\cdot
    \left(\mathbb 1+\begin{pmatrix}
                        0 \\
                        m^{k+1}
                    \end{pmatrix}e^T_{k+1}\right)                                                     \\
     & =\mathbb 1+\begin{pmatrix}
                      0 \\
                      m^{k+1}
                  \end{pmatrix}e^T_{k+1}+
    \begin{pmatrix}
        0 \\
        m^k
    \end{pmatrix}e^T_k+\underbrace{\begin{pmatrix}
                                           0 \\
                                           m^k
                                       \end{pmatrix}e^T_k
    \begin{pmatrix}
            0 \\
            m^{k+1}
        \end{pmatrix}e^T_{k+1}}_{\text{Come prima }=0}                                                    \\
     & \operatorname{L}^{k}\cdot\operatorname{L}^{k+1}=\mathbb 1+\begin{pmatrix}
                                                                     0 \\
                                                                     m^k
                                                                 \end{pmatrix}e^T_k+
    \begin{pmatrix}
        0 \\
        m^{k+1}
    \end{pmatrix}e^T_{k+1}
\end{align*}
Volevamo calcolare:
\begin{align*}
     & \left(\widetilde{\operatorname{L}}^1\right)^{-1}\left(\widetilde{\operatorname{L}}^2\right)^{-1}\dots\left(\widetilde{\operatorname{L}}^{(n-1)}\right)^{-1}=\operatorname{L}^1\operatorname{L}^2\dots\operatorname{L}^{n-1} \\
     & =\mathbb 1+\begin{pmatrix}
                      0 \\
                      m^{1}
                  \end{pmatrix}e^T_{1}+\begin{pmatrix}
                                           0 \\
                                           m^{2}
                                       \end{pmatrix}e^T_{2}\dots\begin{pmatrix}
                                                                    0 \\
                                                                    m^{n-1}
                                                                \end{pmatrix}e^T_{n-1}
\end{align*}

Se vale l'\textbf{ipotesi} iniziale ad ogni passo e scelgo gli $m$ secondo $\circledast$ ottengo $A=LU$, come volevamo $\checkmark$.

\colorbox{yellow}{Osservazione}: Se $A^{(k)}(k,k)$ detto \textbf{Pivot}, è molto piccolo in modulo. La scelta $\circledast$, causa
moltiplicatori molto grandi e le sottrazioni per eliminare gli elementi sottodiagonali, diventano fortemente \textbf{instabili}.

\colorbox{yellow}{Conclusione}: \underline{$LU$ così non può essere calcolato}.

\colorbox{yellow}{Osservazione}: La limitazione che deriva dall'\textbf{ipotesi} si può rimuovere scambiando
la $k$-esima riga con quella $l$-esima con $l>k$.
$\exists$ sicuramente $l>k:\ A^{(K)}(l,k)\ne0$.

\subsection{Matrici di Permutazione}
\begin{defi}{Matrice di Permutazione}{}
    Sia $P\in M_{n\times n}(\mathbb R):\ P_{ij}=\begin{cases}
            1 \\0
        \end{cases}$ Tale per cui esiste una permutazione $\sigma$ dell'insieme $\{1,2,\dots,n\}$, può essere associata alla matrice $$P_{ij}:=\begin{cases}
            1\text{ se } j=\sigma(i) \\
            0\text{ altrimenti}
        \end{cases}$$
    Ossia la matrice che si ottiene permutando con $\sigma$ le righe di $\mathbb 1_n$.
\end{defi}
Le matrici di permutazione hanno delle proprietà notevoli:
\begin{enumerate}
    \item Se $P,Q$ sono matrici di permutazione $\Rightarrow PQ,QP$ sono di permutazione.
    \item $PP^T=P^TP=\mathbb 1_n$ sono \textbf{ortogonali}.
    \item $\det P=\pm 1$ quelle con determinante positivo/negativo sono dette \textbf{pari}/\textbf{dispari}.
    \item Se $\sigma$ è lo scambio di due elementi, allora $P$ è detta di scambio e $P=P^T=P^{-1}$
\end{enumerate}

Se $P$ è una matrice di permutazione che rappresenta $\sigma$ ed $A$ ha dimensione uguale a $P$, allora
\begin{itemize}
    \item $AP_\sigma=A$ (moltiplicazione a destra) significa permutare le colonne di $A$, ossia $\left(AP_\sigma\right)_{:,j}=A_{:,\sigma(j)}$
    \item $P_\sigma A=A$ (moltiplicazione a sinistra) significa permutare le righe di $A$, ossia $\left(P_\sigma A\right)_{i,:}=A_{\sigma(i),:}$

\end{itemize}

Le permutazioni più semplici sono gli scambi, se $\sigma$ è uno scambio $\left(k\to l,l\to k\right)$, allora
vuole dire che $P_\sigma$ ha la diagonale con tutti $1$ e alla riga $k$ colonna $l$ ho $1$ e anche alla riga $l$ colonna $k$, mentre il resto rimane uguale.\\
Equivale a prendere $\mathbb 1$ ed effettuare lo scambio o su righe o su colonne.

Se si aggiunge nel vecchio algoritmo la possibilità di scambiare le righe, ossia il \textbf{pivoting parziale} $\times$ righe; si passa da:
\begin{align*}
    &A^1=A\\
    &A^k=\widetilde{\operatorname{L}}^{(k-1)}A^{(k-1)}\quad k=2\dots n\\
\end{align*}
e diventa:
\begin{align*}
    &\widetilde{A}^1=A\\
    &A^1=P^1A\\
    &\widetilde{A}^k=\widetilde{L}^{(k-1)}A^{k-1}\\
    &A^k=P^k\widetilde{A}^k
\end{align*}
Ossia ad ogni passo permuto e mi calcolo i moltiplicatori. Successivamente calcolo $\widetilde{L}^{(k-1)}$, poi moltiplico.

\begin{defi}{Pivoting}{}
    Fattorizzazione $LU$ con \underline{pivoting} parziale per righe:
    \[\boxed{LU=PA}\]
    Quindi $Ax=b\Leftrightarrow PAx=Pb$ e $LUx=Pb$.
\end{defi}
Lo scambio permette di avere la possibilità di calcolo. La strategia di scelta del \textbf{pivot} ci consente di non averne uno troppo piccolo, che nel caso non andrebbe bene 
dato che ci uscirebbe un moltiplicatore grande e avremmo una sottrazione tra due numeri grandi.

Scelgo, quindi, $l>k\ :\ \left|A(l,k)\right|=\max_{p>k}\left|A(p,k)\right|$, ossia scelgo la riga in modo che il suo elemento nella posizione $k$ sia massimo in modulo rispetto agli elementi
sotto la diagonale in una certa colonna.

\[P^k=\mathbb 1_{k-1}\]

Risolviamo $Ly=Pb$ con $Ux=y$ usando la sostituzione in avanti e all'indietro.

La Fattorizzazione ha costo totale:\[\sum_{k=1}^{n-1}\left[(n-k)+2(n-k)^2\right]=\boxed{0(n^3)}\]
Di cui $(n-k)$ è dato dalle divisioni, mentre $2(n-k)^2$ è dato dalle moltiplicazioni $+$ le somme.

\colorbox{yellow}{Osservazione}: In moltissime applicazioni le matrici sono \underline{sparse}, ossia molti elementi sono nulli.
\begin{note}{Matrici Sparse}{}
    Una successione di matrici $A^{(n)}\in M_{n\times n}(\mathbb R)$ è \textbf{sparsa} se $n\in\left(A^{(n)}\right)=O(n)$
\end{note}
\end{document}